\chapter{Introdução}
\label{intro}

% deve substituir as duas linhas seguintes pelo texto da introdução

Ao longo da licenciatura em Engenharia Informática, manifestei bastante curiosidade em compreender e aprender mais sobre várias áreas, mas duas destacaram-se: Engenharia de Software e Sistemas de Informação. No terceiro ano, surgiu a necessidade de concluir a licenciatura com um estágio ou um projeto. Optei por fazer um projeto relacionado com as duas unidades curriculares anteriormente mencionadas. Optei por um projeto de análise de dados, com a orientação da docente Isabel Sofia Brito.

Este projeto teve início através de um e-mail, com um pedido de mentoria para a execução do projeto por parte do docente, e de seguida com a aprovação da comissão do curso. Após a aprovação destes passos, iniciámos o desenvolvimento do projeto. No primeiro contacto com o cliente observamos então o problema mencionado na proposta do projeto, a qualidade do ar. 

A qualidade do ar constitui um dos principais indicadores ambientais relacionados 
com a saúde pública, a sustentabilidade e a qualidade de vida das populações. O 
aumento das emissões provenientes do tráfego rodoviário, atividades industriais, 
agricultura e fenómenos naturais torna indispensável a monitorização contínua dos 
poluentes atmosféricos, permitindo apoiar decisões fundamentadas por parte das 
entidades públicas e sensibilizar os cidadãos.
Em Portugal, a avaliação e gestão da qualidade do ar encontram-se enquadradas 
pelo Decreto-Lei n.º 102/2010, que estabelece os objetivos da qualidade do ar 
ambiente, define metodologias de avaliação, requisitos de qualidade dos dados e 
determina que a informação relativa à qualidade do ar seja disponibilizada ao 
público. Entre os seus objetivos destacam-se a prevenção e redução dos efeitos 
nocivos para a saúde humana e para o ambiente, a avaliação da qualidade do ar 
através de métodos normalizados e a disponibilização de informação fiável e atualizada aos cidadãos.
No primeiro contacto com a cliente compreendermos o problema, que se situa em Beja, os dados que usamos como dataset, tem origem em quatro estações meteorológicas/sensores que estão situados em Beja, nos  seguintes pontos IPBeja, OviBeja, Parque de trabalhos (rotunda força aérea - burger king), Jardim Publico. 
